\appendices
\section{Tabela de contribuição da equipe}
\label{ap:contribuicao}

Conforme a premissa 13 do roteiro, este apêndice não conta no limite de páginas. As atividades foram levantadas dos históricos de versão da GitPage e do histórico de \textit{commits}. O percentual de esforço deve ser preenchido pela equipe a partir da planilha de esforço e totalizar 100\%.

\begin{table}[!h]
	\caption{Contribuição de cada membro na Entrega 1}
	\label{tab:contribuicao}
	\centering
	\footnotesize
	\begin{tabularx}{\columnwidth}{L{1.9cm} Y C{0.9cm}}
		\toprule
		\textbf{Membro} & \textbf{Atividades} & \textbf{Esforço} \\
		\midrule
		Caio Melo Borges & Fase 1 completa (propósito, partes interessadas, software avaliado, modelo de qualidade, matriz Impacto~$\times$~Risco, escopo, ODS); síntese da Fase 1; referências da Fase 1 & \pendente{--} \\
		Daniel Nunes Duarte & \pendente{Descrever as atividades (sem registro no repositório até 19/10/2026)} & \pendente{--} \\
		Davi dos Santos Brito Nobre & Repositório, estrutura MkDocs e GitPage; primeira versão da Fase 1; integração do PR da Fase 1 & \pendente{--} \\
		Felipe de Aquino Campêlo & Fase 2: visão geral, objetivos G1 a G4, questões Q1 a Q14, métricas M1 a M19, níveis de pontuação e gráficos GQM & \pendente{--} \\
		Leonardo Fachinello Bonetti & Fase 2: M20, MA1, hipóteses H1 a H14 e HR1 a HR4, critérios de julgamento, rastreabilidade, referências; integração do PR da Fase 2 & \pendente{--} \\
		\pendente{Responsável a definir} & Artigo no modelo IEEE, questionário com estudantes e gestão PSM/CID (\textit{commits} de arthurlleite) & \pendente{--} \\
		\midrule
		\textbf{Total} & & \textbf{100\%} \\
		\bottomrule
	\end{tabularx}
	\fonte{elaborado pela equipe.}
\end{table}
